% Options for packages loaded elsewhere
\PassOptionsToPackage{unicode}{hyperref}
\PassOptionsToPackage{hyphens}{url}
%
\documentclass[
]{article}
\usepackage{amsmath,amssymb}
\usepackage{iftex}
\ifPDFTeX
  \usepackage[T1]{fontenc}
  \usepackage[utf8]{inputenc}
  \usepackage{textcomp} % provide euro and other symbols
\else % if luatex or xetex
  \usepackage{unicode-math} % this also loads fontspec
  \defaultfontfeatures{Scale=MatchLowercase}
  \defaultfontfeatures[\rmfamily]{Ligatures=TeX,Scale=1}
\fi
\usepackage{lmodern}
\ifPDFTeX\else
  % xetex/luatex font selection
\fi
% Use upquote if available, for straight quotes in verbatim environments
\IfFileExists{upquote.sty}{\usepackage{upquote}}{}
\IfFileExists{microtype.sty}{% use microtype if available
  \usepackage[]{microtype}
  \UseMicrotypeSet[protrusion]{basicmath} % disable protrusion for tt fonts
}{}
\makeatletter
\@ifundefined{KOMAClassName}{% if non-KOMA class
  \IfFileExists{parskip.sty}{%
    \usepackage{parskip}
  }{% else
    \setlength{\parindent}{0pt}
    \setlength{\parskip}{6pt plus 2pt minus 1pt}}
}{% if KOMA class
  \KOMAoptions{parskip=half}}
\makeatother
\usepackage{xcolor}
\usepackage[margin=1in]{geometry}
\usepackage{color}
\usepackage{fancyvrb}
\newcommand{\VerbBar}{|}
\newcommand{\VERB}{\Verb[commandchars=\\\{\}]}
\DefineVerbatimEnvironment{Highlighting}{Verbatim}{commandchars=\\\{\}}
% Add ',fontsize=\small' for more characters per line
\usepackage{framed}
\definecolor{shadecolor}{RGB}{248,248,248}
\newenvironment{Shaded}{\begin{snugshade}}{\end{snugshade}}
\newcommand{\AlertTok}[1]{\textcolor[rgb]{0.94,0.16,0.16}{#1}}
\newcommand{\AnnotationTok}[1]{\textcolor[rgb]{0.56,0.35,0.01}{\textbf{\textit{#1}}}}
\newcommand{\AttributeTok}[1]{\textcolor[rgb]{0.13,0.29,0.53}{#1}}
\newcommand{\BaseNTok}[1]{\textcolor[rgb]{0.00,0.00,0.81}{#1}}
\newcommand{\BuiltInTok}[1]{#1}
\newcommand{\CharTok}[1]{\textcolor[rgb]{0.31,0.60,0.02}{#1}}
\newcommand{\CommentTok}[1]{\textcolor[rgb]{0.56,0.35,0.01}{\textit{#1}}}
\newcommand{\CommentVarTok}[1]{\textcolor[rgb]{0.56,0.35,0.01}{\textbf{\textit{#1}}}}
\newcommand{\ConstantTok}[1]{\textcolor[rgb]{0.56,0.35,0.01}{#1}}
\newcommand{\ControlFlowTok}[1]{\textcolor[rgb]{0.13,0.29,0.53}{\textbf{#1}}}
\newcommand{\DataTypeTok}[1]{\textcolor[rgb]{0.13,0.29,0.53}{#1}}
\newcommand{\DecValTok}[1]{\textcolor[rgb]{0.00,0.00,0.81}{#1}}
\newcommand{\DocumentationTok}[1]{\textcolor[rgb]{0.56,0.35,0.01}{\textbf{\textit{#1}}}}
\newcommand{\ErrorTok}[1]{\textcolor[rgb]{0.64,0.00,0.00}{\textbf{#1}}}
\newcommand{\ExtensionTok}[1]{#1}
\newcommand{\FloatTok}[1]{\textcolor[rgb]{0.00,0.00,0.81}{#1}}
\newcommand{\FunctionTok}[1]{\textcolor[rgb]{0.13,0.29,0.53}{\textbf{#1}}}
\newcommand{\ImportTok}[1]{#1}
\newcommand{\InformationTok}[1]{\textcolor[rgb]{0.56,0.35,0.01}{\textbf{\textit{#1}}}}
\newcommand{\KeywordTok}[1]{\textcolor[rgb]{0.13,0.29,0.53}{\textbf{#1}}}
\newcommand{\NormalTok}[1]{#1}
\newcommand{\OperatorTok}[1]{\textcolor[rgb]{0.81,0.36,0.00}{\textbf{#1}}}
\newcommand{\OtherTok}[1]{\textcolor[rgb]{0.56,0.35,0.01}{#1}}
\newcommand{\PreprocessorTok}[1]{\textcolor[rgb]{0.56,0.35,0.01}{\textit{#1}}}
\newcommand{\RegionMarkerTok}[1]{#1}
\newcommand{\SpecialCharTok}[1]{\textcolor[rgb]{0.81,0.36,0.00}{\textbf{#1}}}
\newcommand{\SpecialStringTok}[1]{\textcolor[rgb]{0.31,0.60,0.02}{#1}}
\newcommand{\StringTok}[1]{\textcolor[rgb]{0.31,0.60,0.02}{#1}}
\newcommand{\VariableTok}[1]{\textcolor[rgb]{0.00,0.00,0.00}{#1}}
\newcommand{\VerbatimStringTok}[1]{\textcolor[rgb]{0.31,0.60,0.02}{#1}}
\newcommand{\WarningTok}[1]{\textcolor[rgb]{0.56,0.35,0.01}{\textbf{\textit{#1}}}}
\usepackage{graphicx}
\makeatletter
\def\maxwidth{\ifdim\Gin@nat@width>\linewidth\linewidth\else\Gin@nat@width\fi}
\def\maxheight{\ifdim\Gin@nat@height>\textheight\textheight\else\Gin@nat@height\fi}
\makeatother
% Scale images if necessary, so that they will not overflow the page
% margins by default, and it is still possible to overwrite the defaults
% using explicit options in \includegraphics[width, height, ...]{}
\setkeys{Gin}{width=\maxwidth,height=\maxheight,keepaspectratio}
% Set default figure placement to htbp
\makeatletter
\def\fps@figure{htbp}
\makeatother
\setlength{\emergencystretch}{3em} % prevent overfull lines
\providecommand{\tightlist}{%
  \setlength{\itemsep}{0pt}\setlength{\parskip}{0pt}}
\setcounter{secnumdepth}{-\maxdimen} % remove section numbering
\ifLuaTeX
  \usepackage{selnolig}  % disable illegal ligatures
\fi
\IfFileExists{bookmark.sty}{\usepackage{bookmark}}{\usepackage{hyperref}}
\IfFileExists{xurl.sty}{\usepackage{xurl}}{} % add URL line breaks if available
\urlstyle{same}
\hypersetup{
  hidelinks,
  pdfcreator={LaTeX via pandoc}}

\author{}
\date{\vspace{-2.5em}}

\begin{document}

\hypertarget{benchmarking-memo-3-team-follow-up-spatial-variation-and-a-sign-error-in-the-benchmarking-formula}{%
\section{Benchmarking Memo 3: Team Follow-Up -- Spatial Variation and a
Sign Error in the Benchmarking
Formula}\label{benchmarking-memo-3-team-follow-up-spatial-variation-and-a-sign-error-in-the-benchmarking-formula}}

\textbf{Date:} July 2, 2026 \textbf{Prepared by:} PM, with assistance
from Claude Code Sonnet 5 \textbf{For:} The full SNVDEM-COL team,
following the 2026-07-02 meeting

\begin{center}\rule{0.5\linewidth}{0.5pt}\end{center}

\hypertarget{how-this-memo-fits-with-what-came-before}{%
\subsection{How this memo fits with what came
before}\label{how-this-memo-fits-with-what-came-before}}

This is the third memo in a chain, and each one only makes sense in
light of the last:

\begin{enumerate}
\def\labelenumi{\arabic{enumi}.}
\tightlist
\item
  \textbf{Memo 1} (\texttt{benchmarking\_memo.md}, June 26) laid out the
  problem: our two dimensions -- \texttt{EL\_col\_mt} (elections,
  V-Dem's latent scale) and \texttt{CL\_col\_mt} (civil liberties,
  {[}0,1{]} scale) -- couldn't be combined into a single
  \texttt{sndem\_final} without distorting or flattening the trend.
  Every method tried produced a bad result.
\item
  \textbf{MC's email reply} (July 2) resolved the conceptual question:
  benchmarked scores are the intended final measure for all uses, and
  the two dimensions should be combined by converting
  \texttt{CL\_col\_mt} to a Z-score with \texttt{qnorm()}.
\item
  \textbf{Memo 2} (\texttt{benchmarking\_memo2\_2026-07-02.md})
  implemented that, caught and fixed a mistake in the first attempt (an
  ipsative -- Colombia-only -- standardization that would have broken
  comparability with other countries), and validated the result against
  the raw V-Dem national anchors.
\item
  \textbf{This memo} follows up on the 2026-07-02 team meeting, where
  Patrick presented the Memo 2 maps. It answers Patrick's own follow-up
  question about how the Z-scores are calculated, and reports something
  found while digging into why the benchmarked maps show less visible
  spatial variation than the unbenchmarked ones -- a sign error in the
  benchmarking formula, present since the original scripts and unrelated
  to anything in Memos 1-2.
\end{enumerate}

\textbf{If you only read one section, read Section 2.}

\begin{center}\rule{0.5\linewidth}{0.5pt}\end{center}

\hypertarget{recap-what-benchmarked-and-unbenchmarked-each-show}{%
\subsection{1. Recap: what ``benchmarked'' and ``unbenchmarked'' each
show}\label{recap-what-benchmarked-and-unbenchmarked-each-show}}

For readers who weren't in the weeds of Memos 1-2:

\begin{itemize}
\item
  \textbf{Unbenchmarked (\texttt{sndem}, \texttt{snelect},
  \texttt{sncivlib}):} each municipality's geographic/structural
  predictors (rurality, distance from capital, violence, etc.), weighted
  by V-Dem expert judgments about what matters for elections and civil
  liberties, then standardized against the empirical distribution of
  \emph{other Colombian municipalities} (a CDF, 2000-2023 pooled). This
  is a \textbf{within-Colombia, relative} measure -- it tells you how a
  municipality compares to the rest of Colombia, but not how Colombia as
  a whole compares to any other country, and only weakly reflects
  absolute national-level change over time (Memo 1, Sections 3-4).
\item
  \textbf{Benchmarked (\texttt{EL\_col\_mt}, \texttt{CL\_col\_mt},
  \texttt{sndem\_final}):} each municipality's unbenchmarked score is
  re-expressed as a deviation from the actual V-Dem national estimate
  for that year (\texttt{v2elffelr} for elections,
  \texttt{v2x\_civlib}-derived \texttt{CLSNmean} for civil liberties),
  scaled by how much the country is estimated to vary internally
  (\texttt{weighted\_range}/\texttt{wtdCL\_range}, built from V-Dem's
  own uncertainty and ``unevenness'' codings). This is meant to be a
  measure that's \textbf{anchored to something real and comparable
  across countries} -- which is the entire reason the benchmarking step
  exists (per the operational strategy document quoted in Memo 1: ``we
  conceive of free and fair elections at the municipal level as a set of
  deviations centered around this summary subnational measure'').
\end{itemize}

Michael's point at the meeting -- that both are useful for different
things, and we should keep both -- is right and isn't in question here.
The question this memo investigates is whether the benchmarked version
is currently computing what it's supposed to.

\begin{center}\rule{0.5\linewidth}{0.5pt}\end{center}

\hypertarget{finding-benchmarking-is-currently-inverting-the-spatial-pattern-not-just-muting-it}{%
\subsection{2. Finding: benchmarking is currently inverting the spatial
pattern, not just muting
it}\label{finding-benchmarking-is-currently-inverting-the-spatial-pattern-not-just-muting-it}}

At the meeting, the working explanation for why the new maps show less
visible spatial (municipality-to-municipality) variation than the old
ones was that national-level, year-to-year change is large enough to
visually dominate the color scale, ``overwhelming'' the subtler spatial
differences. That's a real effect and part of the picture -- but
checking it turned up something bigger.

\textbf{Within every single year from 2000-2023, a municipality's
benchmarked score is currently the exact mathematical mirror image of
its unbenchmarked score.} Not ``compressed'' or ``less visible'' --
inverted. The correlation between \texttt{snelect} (unbenchmarked) and
\texttt{EL\_col\_mt} (benchmarked) is \textbf{r = -1.0000} in all 24
years. Same for civil liberties: \texttt{sncivlib}
vs.~\texttt{CL\_col\_mt}, r = -1.0000 in all 24 years. The scatter plot
below (2023, elections) isn't noise or a weak relationship -- it's a
perfectly straight line with negative slope:

\begin{figure}
\centering
\includegraphics{../rank_inversion_scatter_2023.png}
\caption{Rank inversion}
\end{figure}

Concretely: sorting Colombia's municipalities into deciles by their
unbenchmarked elections score, the \emph{top} decile has the
\emph{lowest} mean benchmarked score (1.43) and the \emph{bottom} decile
has the \emph{highest} (2.03) -- a clean, monotonic reversal from decile
1 through decile 10. This isn't specific to any one region or grouping;
it will show up in any comparison between high- and low-scoring
municipalities.

\hypertarget{root-cause}{%
\subsubsection{Root cause}\label{root-cause}}

The benchmarking formula (\texttt{08\_benchmark/01\_benchmark.R},
unchanged from MC's original design):

\begin{verbatim}
EL_col_mt = v2elffelr + (snelect - snelectyrmean) * weighted_range / ELrange_975_025
CL_col_mt = CLSNmean  + (sncivlib - CLSNyrmean)   * wtdCL_range   / CLrange_975_025
\end{verbatim}

\texttt{weighted\_range} and \texttt{wtdCL\_range} are national
constants (same value for every municipality in a given year) meant to
represent \emph{how much} Colombia's subnational units are estimated to
vary -- a magnitude, which should be positive. Checked across the full
panel, \textbf{\texttt{weighted\_range} is negative in all 24 years, and
\texttt{wtdCL\_range} is negative in all 24 years}:

\begin{figure}
\centering
\includegraphics{../weighted_range_negative_by_year.png}
\caption{Negative range multiplier}
\end{figure}

Since \texttt{(snelect\ -\ snelectyrmean)} is positive for above-average
municipalities and negative for below-average ones, multiplying by a
\emph{negative} \texttt{weighted\_range} flips the sign of every
municipality's deviation before it's added back to the national anchor.
That's the entire mechanism -- nothing more exotic than a sign error
propagating through the formula.

\textbf{Why is the multiplier negative?} Tracing it back to
\texttt{06\_vdem\_data/01\_weighting.R} (streamlined rewrite of MC's
original weighting script, same formula preserved unchanged):

\begin{Shaded}
\begin{Highlighting}[]
\NormalTok{weighted\_range }\OtherTok{=}\NormalTok{ ((}\DecValTok{2} \SpecialCharTok{{-}}\NormalTok{ snlsff\_1) }\SpecialCharTok{*}\NormalTok{ HPD }\SpecialCharTok{+}\NormalTok{ (}\DecValTok{2} \SpecialCharTok{{-}}\NormalTok{ snlsff\_0) }\SpecialCharTok{*}\NormalTok{ HPD }\SpecialCharTok{*} \DecValTok{2}\NormalTok{) }\SpecialCharTok{/}\NormalTok{ (snlsff\_0 }\SpecialCharTok{+}\NormalTok{ snlsff\_1 }\SpecialCharTok{+}\NormalTok{ snlsff\_2)}
\end{Highlighting}
\end{Shaded}

\texttt{snlsff\_0}, \texttt{snlsff\_1}, \texttt{snlsff\_2} are
\textbf{counts of V-Dem coders} who selected ``same'' (0), ``somewhat
different'' (1), or ``significantly different'' (2) when asked how
uneven Colombia's subnational elections are. For Colombia in a typical
year, 5+ coders are surveyed, so it's common for \texttt{snlsff\_0} or
\texttt{snlsff\_1} to exceed 2 -- at which point
\texttt{(2\ -\ snlsff\_1)} or \texttt{(2\ -\ snlsff\_0)} goes negative.
Example, 2010: \texttt{snlsff\_0\ =\ 12}, \texttt{snlsff\_1\ =\ 16},
giving \texttt{weighted\_range\ =\ -0.631}. The same mechanic affects
\texttt{wtdCL\_range} via \texttt{clrgunev\_0/1/2}.

The \texttt{(2\ -\ count)} construction only produces a sensible
(positive) result if \texttt{snlsff\_0}/\texttt{snlsff\_1} are bounded
to at most 2 -- i.e., if they were meant to be \textbf{proportions or
weights on a 0-2 scale}, not raw coder counts. The operational strategy
document's own description of the intended design (``subnational level
+/- (\% of coders * range of variation\ldots)'') talks in terms of
\emph{percent of coders}, which supports this reading. Whether the fix
is to convert counts to proportions before this formula, cap the
multiplier at zero, or something else structural is MC's call -- this is
his original formula design, not something introduced in Memos 1-2 or in
Patrick's June/July revisions (confirmed unchanged since
\texttt{weighting\_summer2025.R}).

\hypertarget{what-this-does-and-doesnt-affect}{%
\subsubsection{What this does and doesn't
affect}\label{what-this-does-and-doesnt-affect}}

\begin{itemize}
\tightlist
\item
  \textbf{Affected:} the within-year \emph{spatial} pattern in every
  benchmarked map produced so far, including the maps just presented to
  the team (\texttt{sndem\_final\_facet\_allyears.png},
  \texttt{sndem\_dims\_milestone\_years.png}) --
  municipality-to-municipality comparisons within a given year are
  backwards.
\item
  \textbf{Not affected:} the \emph{national} year-over-year trend.
  \texttt{sndem\_final}'s national mean by year is anchored directly to
  \texttt{v2elffelr}/\texttt{CLSNmean} (validated against the raw V-Dem
  anchors in Memo 2, \texttt{trend\_vs\_national\_anchors.png}) -- that
  validation only checked the national-level series, which doesn't
  depend on the sign of \texttt{weighted\_range}. So the overall
  ``Colombia improved 2000-2018, dipped 2019-2021'' story likely still
  holds; it's specifically the \emph{within-country geography} that's
  currently backwards.
\item
  \textbf{Not affected:} the \texttt{qnorm()}/global-standardization
  work in Memo 2 -- that's a separate, later step in the pipeline and is
  correct on its own terms. It's just operating on
  \texttt{EL\_col\_mt}/\texttt{CL\_col\_mt} inputs that are already
  sign-flipped upstream.
\end{itemize}

\begin{center}\rule{0.5\linewidth}{0.5pt}\end{center}

\hypertarget{answering-patricks-question-from-the-meeting-about-z-scores}{%
\subsection{3. Answering Patrick's question from the meeting about
Z-scores}\label{answering-patricks-question-from-the-meeting-about-z-scores}}

At the meeting, Patrick asked for clarification on how the Z-scores were
calculated and whether they properly represent global benchmarks. Short
answer: yes, as of the fix described in Memo 2 -- \texttt{EL\_col\_gz}
and \texttt{CL\_col\_gz} are standardized against fixed constants from
the full V-Dem country-year panel (2000-2023, all countries), not from
Colombia's own data, so they are comparable across countries in
principle. That part is independent of, and unaffected by, the
sign-error finding in Section 2 above -- the global-standardization fix
operates downstream of \texttt{EL\_col\_mt}/\texttt{CL\_col\_mt} and
would apply equally correctly once those are fixed.

\begin{center}\rule{0.5\linewidth}{0.5pt}\end{center}

\hypertarget{recommendation-before-further-sharing}{%
\subsection{4. Recommendation before further
sharing}\label{recommendation-before-further-sharing}}

\textbf{Hold off on presenting the benchmarked spatial maps further
until the \texttt{weighted\_range}/\texttt{wtdCL\_range} sign issue is
resolved with MC.} The national trend line (the improving-over-time
story) is likely fine to keep using. The within-Colombia geography shown
in the benchmarked maps is currently backwards and shouldn't be
presented as-is. The unbenchmarked maps do not have this problem and
remain valid for showing subnational spatial variation, consistent with
what Michael and Kelly suggested -- presenting both versions with a
clear explanation of what each is for is still the right end state; we
just need the benchmarked version corrected first.

\begin{center}\rule{0.5\linewidth}{0.5pt}\end{center}

\hypertarget{next-steps}{%
\subsection{5. Next steps}\label{next-steps}}

\begin{itemize}
\tightlist
\item[$\square$]
  Confirm the \texttt{weighted\_range}/\texttt{wtdCL\_range} sign issue
  with MC and agree on a fix (proportions vs.~counts, or another
  approach to the formula).
\item[$\square$]
  Re-run \texttt{06\_vdem\_data/01\_weighting.R} -\textgreater{}
  \texttt{07\_weighting/01\_weight\_predictors.R} -\textgreater{}
  \texttt{08\_benchmark/01\_benchmark.R} once fixed, and re-validate
  (rank correlation should be strongly \emph{positive}, not -1, after
  the fix).
\item[$\square$]
  Regenerate the benchmarked maps once corrected.
\item[$\square$]
  Per Michael's suggestion: draft the explanatory language
  distinguishing benchmarked (cross-national, smoother, anchored to
  V-Dem) from unbenchmarked (within-Colombia, more spatial detail) for
  the paper, once the benchmarked maps are corrected.
\item[$\square$]
  Revisit Memo 2's six open questions (reference window, reference
  population, pooled vs.~year-specific global constants,
  \texttt{qnorm()} order of operations, etc.) -- those remain relevant
  once this more urgent issue is resolved.
\item[$\square$]
  Explore how the corrected spatial-variation data could support
  cross-national comparison hypotheses (per Michael's example of testing
  democratic influence across national borders) once the sign issue is
  fixed.
\end{itemize}

\begin{center}\rule{0.5\linewidth}{0.5pt}\end{center}

\hypertarget{files}{%
\subsection{Files}\label{files}}

\begin{itemize}
\tightlist
\item
  \texttt{08\_benchmark/04\_spatial\_rank\_check.R} -- this
  investigation (rank-inversion check, plots)
\item
  \texttt{08\_benchmark/weighted\_range\_negative\_by\_year.png},
  \texttt{rank\_inversion\_scatter\_2023.png}
\item
  \texttt{08\_benchmark/memo/benchmarking\_memo.md} -- Memo 1
\item
  \texttt{08\_benchmark/memo/benchmarking\_memo2\_2026-07-02.md} -- Memo
  2
\item
  \texttt{HANDOFF\_pipeline\_revision\_june2026.md} -- running pipeline
  log, to be updated once this is resolved
\end{itemize}

\end{document}
